\chapter{External literature results}
\label{chap:external-results}

This flat chapter contains only precise published theorem statements that the project consumes without reproving.  Every node is stated using objects already introduced above and retains its source and locator.

\section{Published differentials, extensions, and stable-operation results}

% SOURCE: aimpaper/main.tex:2493-2567 and cited Toda/Moss identities; explicit load-bearing identities.
\begin{theorem}[Toda product and shuffle identities]
  \label{thm:toda-product-identities}
  \uses{def:toda-bracket,def:massey-product,def:adams-detection,def:external-result}
  \notready
  Explicitly located stable-homotopy results supply the Toda naturality,
  suspension, product, and shuffle containments used in Section 7, with their
  full indeterminacy.  In the particular degrees and chosen nullhomotopies of
  the paper they supply
  \[
    \eta^2\in\langle[h_0],\eta,[h_0]\rangle
  \]
  and, for every
  $b\in\langle\lambda^3\alpha_1,[h_0],\eta\rangle$,
  \[
    b[h_0]=\eta\lambda^3\eta\alpha_1.
  \]
  The record also carries the zero-indeterminacy checks needed for these
  equalities; a general containment is never silently strengthened to an
  equality.
\end{theorem}

\begin{proof}
  Instantiate the standard Toda composition and shuffle theorems with the
  paper's chosen nullhomotopies, calculate both indeterminacy subgroups, and
  use the typed vanishing evidence in the two relevant stems to show they are
  zero.  Transport the Hopf relation and signs to the characteristic-two
  convention.
\end{proof}

% SOURCE: aimpaper/main.tex:140--150, specialized to j=4.
% VERIFIED-IN: KIP126/Classical/Adams/Basic.lean, declaration
% KIP126.Classical.Adams.cataloguedAdamsOneLine_h₄_degrees_bound.
\begin{theorem}[$h_4$ $d_2$ tracer bullet]
  \label{thm:external-adams-h4-d2-slice}
  \uses{def:classical-adams-e2-e3-shape-slice,def:catalogued-external-result}
  \lean{KIP126.Classical.Adams.cataloguedAdamsOneLine_h₄_degrees_bound}
  \leanok
  For a chosen spectrum-bound sphere $E_2$ presentation with its lawful
  bilinear, unital, associative, Leibniz-compatible and nondegenerate
  internal multiplication, the located \texttt{adams\_one\_line} result is
  consumed together with the convergence-and-detection witnesses and supplies the
  relation
  \[
    d_2(h_4)=h_0h_3^2.
  \]
  The source is in bidegree $(1,16)$ and the target in $(3,17)$.  The consumer
  is tied to the canonical claim root, owner, trust class, and locator in the
  external claim ledger; it does not turn the literature result into a project
  axiom.
\end{theorem}

\begin{proof}
  The exact consumer extracts the located result, retains the chosen
  spectrum's convergence-and-detection data, checks all lawful algebra fields, and
  derives the displayed bidegrees from the actual Mathlib $E_2$ differential
  component and the supplied sphere product.
\end{proof}

% SOURCE: aimpaper/main.tex:140-150; citations Adams/Mahowald--Tangora/May.
% VERIFIED-IN: aimpaper/main.tex:140-150 contains a typo: survival is j <= 3, while d_2 is nonzero for j >= 4.
\begin{theorem}[Adams one-line differentials]
  \label{thm:external-adams-one-line}
  \uses{def:adams-hj,def:external-result}
  \notready
  An explicit external result supplies that $h_0,h_1,h_2,h_3$ survive and
  that, for $j\ge4$,
  \[
    d_2(h_j)=h_0h_{j-1}^2\ne0.
  \]
  This is the mathematically consistent reading of the two consecutive
  sentences in \texttt{aimpaper/main.tex:140--150}; the printed ``$j\ge3$
  survives'' is recorded as a source typo, not adopted as a project theorem.
\end{theorem}

\begin{proof}
  This node is a literature input.  Instantiate the external-result record at
  the cited theorem and transport its grading to
  Definition~\ref{def:adams-bidegree}; no internal proof of the classical
  calculation is claimed.
\end{proof}

% SOURCE: aimpaper/main.tex:253-263; bibliography key Adams58.  The local
% source inventory does not contain a checked copy of the original convergence
% proof, so this remains a provenance-carrying input rather than a ready node.
\begin{theorem}[Classical Adams convergence and detection input]
  \label{thm:external-classical-adams-convergence}
  \uses{def:admissible-spectrum,def:ss-convergence-detection-witness,def:external-result}
  \notready
  For a 2-complete connective spectrum of finite type, an explicit external
  result supplies a convergence-and-detection witness for the mod--$2$ Adams
  spectral sequence: its algebraic $E_\infty$ page is the associated graded of
  the 2-complete stable homotopy filtration, and its nonzero permanent classes
  detect the corresponding filtered homotopy elements.  The project records
  this as the standing hypothesis used in
  \texttt{aimpaper/main.tex:253--263}, with the classical source cited there
  (bibliography key \texttt{Adams58}); it does not prove convergence for
  arbitrary unbounded filtered complexes.
\end{theorem}

\begin{proof}
  Instantiate the source-bearing external-result record at the admissible
  spectrum and transport its page and filtration conventions to the concrete
  classical Adams sequence.  The source locator and the witness remain part of
  the hypothesis consumed by later definitions.
\end{proof}

% SOURCE: Ravenel, Complex Cobordism, Theorem 2.2.14; cited at aimpaper/main.tex:273-275.
\begin{theorem}[Factorization criterion for map filtration]
  \label{thm:external-map-filtration-factorization}
  \uses{def:adams-filtration,def:hf2-homology,def:external-result}
  \notready
  An explicit literature result supplies the implication used by the paper:
  if $\operatorname{AF}(f)\ge k$, then $f$ is represented by a composite of
  $k$ maps, each inducing zero on $H\mathbb F_2$-homology.  The converse is
  included only to the extent stated by the cited theorem and the chosen
  Adams-tower model.
\end{theorem}

\begin{proof}
  Instantiate the source-bearing external result at the chosen Adams tower and
  transport its filtration convention to Definition~\ref{def:adams-filtration}.
  No factorization theorem is postulated globally without this parameter.
\end{proof}

% SOURCE: May, The Additivity of Traces in Triangulated Categories, Lemma 4.6 and TC3; aimpaper/main.tex:1755-1778.
\begin{theorem}[May's smash-boundary lemma]
  \label{thm:external-may-smash-boundary}
  \uses{def:smash-triangle,def:external-result}
  \notready
  Given two distinguished triangles, form their $3\times3$ smash diagram.  If
  $a\in\pi_n(X\wedge Z')$ and $b\in\pi_n(Y\wedge Y')$ have the same image in
  $\pi_n(Y\wedge Z')$, then there is
  $c\in\pi_n(Z\wedge X')$ whose image agrees with $b$ in
  $\pi_n(Z\wedge Y')$ and whose boundary agrees with the boundary of $a$ in
  $\pi_{n-1}(X\wedge X')$.
\end{theorem}

\begin{proof}
  This is supplied as an external result located at May's Lemma 4.6 and axiom
  TC3.  The cited pullback object lifts the compatible pair $(a,b)$; its third
  structure map gives $c$, and commutativity supplies the two equalities.
\end{proof}

\section{Published synthetic and comparison results}

% SOURCE: Pstragowski, Synthetic Spectra and the Cellular Motivic Category, category construction and synthetic analogue functor; summarized at aimpaper/main.tex:755-758.
\begin{theorem}[Existence and functoriality of the synthetic context]
  \label{thm:external-synthetic-foundation}
  \uses{def:synthetic-context,def:external-result}
  \notready
  An explicit Pstragowski external result supplies the synthetic category and
  functor of Definition~\ref{def:synthetic-context}, including stability,
  presentability, the symmetric monoidal structure, lax monoidality of $\nu$,
  and preservation of filtered colimits.
\end{theorem}

\begin{proof}
  Store the cited construction and the synthetic-analogue functor as a
  source-bearing external result, then package only the listed operations and
  coherences into the project interface.  No claim that $\nu$ preserves all
  cofibers is included.
\end{proof}

% SOURCE: Pstragowski, Corollary 4.45
% (label cor:ctau_is_an_algebra; reference/Pst/source/synthetic_spectra.tex:2247-2253).
% VERIFIED-IN: BHSmot, Appendix B, Lemma prop:mod-fil
% (reference/BHSmot/source/Filtered.tex:95-105).
\begin{theorem}[$\lambda$-quotient ring structure]
  \label{thm:external-lambda-quotient-ring}
  \uses{def:synthetic-lambda-quotient,def:external-result}
  \notready
  The cofiber $S^{0,0}/\lambda$ has a unique commutative, equivalently
  $E_\infty$, algebra structure compatible with the canonical quotient map
  $S^{0,0}\to S^{0,0}/\lambda$.
\end{theorem}

\begin{proof}
  Instantiate Pstragowski's located corollary; the filtered-object construction
  in BHSmot gives an independent checked copy for $C\tau$.  This external
  theorem concerns the first cofiber only and supplies no assertion about
  $S^{0,0}/\lambda^n$ for $n>1$.
\end{proof}

% SOURCE: Pstragowski, theorem
% thm:tau_invertible_synthetic_spectra_are_just_spectra
% (reference/Pst/source/synthetic_spectra.tex:2355-2360), together with
% Proposition prop:spectral_yoneda_embedding_the_tau_inversion_of_the_synthetic_analogue
% (reference/Pst/source/synthetic_spectra.tex:2344-2346).
% VERIFIED-IN: reference/BHS/source/SynRevBigraded.tex:34-40.
\begin{theorem}[$\lambda$-inversion recovers classical spectra]
  \label{thm:external-lambda-inversion}
  \uses{def:synthetic-context,def:synthetic-lambda-quotient,def:external-result}
  \notready
  The spectral Yoneda embedding identifies the category of all spectra
  $\operatorname{Sp}$ with the full subcategory of $\lambda$-invertible
  $H\mathbb F_2$-synthetic spectra, and this equivalence is canonically
  symmetric monoidal.  For every spectrum $X$, the canonical map from $\nu X$
  to its spectral Yoneda image is a $\lambda$-inversion; hence
  $\lambda^{-1}\nu X\simeq X$ naturally.  No completion restriction is part of
  this generic-fiber statement.
\end{theorem}

\begin{proof}
  Apply the two located Pstragowski results and retain their monoidal and
  naturality data.  The separate restriction to hypercomplete synthetic
  spectra and $H\mathbb F_2$-local spectra occurs only later in Pstragowski
  (\texttt{reference/Pst/source/synthetic\_spectra.tex:2867--2874}) and is not
  used here.  Transport of a later map or triangle must cite this node and its
  relevant compatibility explicitly.
\end{proof}

% SOURCE: Pstragowski Lemma 4.23; aimpaper/main.tex:759-767, Proposition prop:1f7950df.
\begin{theorem}[Cofiber criterion for $\nu$]
  \label{thm:external-nu-cofiber-criterion}
  \uses{def:synthetic-context,def:cofiber-triangle,def:hf2-homology,def:external-result}
  \notready
  For a cofiber sequence $X\to Y\to Z$, the sequence
  $\nu X\to\nu Y\to\nu Z$ is a synthetic cofiber sequence if and only if
  \[
    0\to H_*(X;\mathbb F_2)\to H_*(Y;\mathbb F_2)
      \to H_*(Z;\mathbb F_2)\to0
  \]
  is short exact.
\end{theorem}

\begin{proof}
  This is the specialization of Pstragowski's Lemma 4.23 carried by an
  external-result parameter.  Check that the project's homology and cofiber
  conventions match the cited statement, then transport it into the abstract
  contexts.
\end{proof}

% SOURCE: BHS Theorem A.8; aimpaper/main.tex:799-807, Theorem thm:rigid.
\begin{theorem}[Synthetic Adams rigidity]
  \label{thm:external-synthetic-rigidity}
  \uses{def:classical-adams-ss,def:synthetic-adams-ss,def:bhs-aim-regrading,def:synthetic-lambda-quotient,def:external-result}
  \notready
  For every admissible $X$, a BHS external result gives
  \[
    {}^{\mathrm{syn}}E_2^{*,*,*}(\nu X)
      \cong E_2^{*,*}(X)\otimes\mathbb F_2[\lambda].
  \]
  Every classical differential $d_r(x)=y$ corresponds to the synthetic
  differential $d_r(x)=\lambda^{r-1}y$, and all synthetic Adams differentials
  arise this way, after applying the regrading adapter.
\end{theorem}

\begin{proof}
  Instantiate BHS Theorem A.8, transport it along
  Definition~\ref{def:bhs-aim-regrading}, and verify the three coordinates of
  each differential.  The resulting isomorphism and differential
  correspondence remain an explicit external-result parameter.
\end{proof}

% SOURCE: BHS Theorem A.1; aimpaper/main.tex:813-823, Theorem thm:17e90ac0.
% VERIFIED-IN: aimpaper/main.tex:819 says tau-Bockstein; in the HF2-synthetic notation it is lambda-Bockstein.
\begin{theorem}[$\lambda$-Bockstein comparison]
  \label{thm:external-lambda-bockstein}
  \uses{thm:external-synthetic-rigidity,def:synthetic-lambda-quotient,def:external-result}
  \notready
  The synthetic Adams spectral sequence of $\nu X$ is isomorphic to the
  $\lambda$-Bockstein spectral sequence.  In particular,
  $E_2^{s,t}(X)\cong\pi_{t-s,t}(\nu X/\lambda)$, and a classical
  $d_r(x)=y$ is represented by a lift of $x$ to
  $\nu X/\lambda^{r-1}$ whose Bockstein boundary is $y$.  The ``$\tau$'' in
  \texttt{aimpaper/main.tex:819} is corrected to $\lambda$ in this context.
\end{theorem}

\begin{proof}
  Apply BHS Theorem A.1 over $\mathbb F_2$, where the sign ambiguity vanishes,
  and transport the grading through the AIM adapter.  Match the cofiber
  boundary with the project's $\lambda$-quotient convention.
\end{proof}

% SOURCE: BHS Corollary A.9; aimpaper/main.tex:841-847, Proposition prop:30e8b746.
\begin{theorem}[$E_\infty$ of $\nu X$]
  \label{thm:external-synthetic-einfty-nu}
  \uses{thm:external-synthetic-rigidity,def:ss-elementwise-page-presentation,def:ss-einfty-presentation,def:external-result}
  \notready
  In AIM grading,
  \[
    E_\infty^{s,t,w}(\nu X)\cong
    \begin{cases}
      Z_\infty^{s,t}(X)/B_{1+t-w}^{s,t}(X),&t\ge w,\\
      0,&t<w.
    \end{cases}
  \]
  This is the regraded form of BHS Corollary A.9.
\end{theorem}

\begin{proof}
  Instantiate the external corollary, apply the BHS-to-AIM regrading, and
  check that the Bockstein filtration is $t-w$.  For nonpositive boundary
  indices use the explicit zero convention, not an implicit natural-number
  subtraction.
\end{proof}

% SOURCE: BHS Corollary A.11; aimpaper/main.tex:849-855, Proposition prop:59f111f.
\begin{theorem}[$E_\infty$ of a finite $\lambda$-quotient]
  \label{thm:external-synthetic-einfty-quotient}
  \uses{thm:external-lambda-bockstein,def:ss-elementwise-page-presentation,def:ss-einfty-presentation,def:external-result}
  \notready
  For $r\ge2$,
  \[
    E_\infty^{s,t,w}(\nu X/\lambda^r)\cong
    \begin{cases}
      Z_{r-t+w}^{s,t}(X)/B_{1+t-w}^{s,t}(X),&0\le t-w<r,\\
      0,&\text{otherwise}.
    \end{cases}
  \]
  This is BHS Corollary A.11 after regrading.
\end{theorem}

\begin{proof}
  Transport the cited result through the grading adapter and identify its
  truncation length with the exponent $r$.  Prove separately that every
  nonpositive cycle or boundary subscript denotes the zero group, so no
  truncated subtraction is hidden in the statement.
\end{proof}

% SOURCE: BHS Lemma 9.15; aimpaper/main.tex:912-919, Proposition prop:ef21f9bc.
\begin{theorem}[Synthetic lift of a filtered map]
  \label{thm:external-synthetic-lift}
  \uses{def:adams-filtration,def:synthetic-context,def:synthetic-lambda-quotient,def:external-result}
  \notready
  If $\operatorname{AF}(f)=k<\infty$, an explicit BHS result supplies a map
  $\widetilde f:\Sigma^{0,k}\nu X\to\nu Y$ with
  $\nu f=\lambda^k\widetilde f$.  This is an arbitrary chosen full lift; the
  result does not make it a component of a lifted distinguished triangle and
  does not assign a cofiber equivalence to it.
\end{theorem}

\begin{proof}
  Instantiate BHS Lemma 9.15 at $f$, transport its suspension convention, and
  retain the chosen factorization and its equality as fields of the external
  result.
\end{proof}

% SOURCE: proof of BHS Lemma 9.15; aimpaper/main.tex:929-940, Proposition prop:41561db2.
\begin{theorem}[Lift of a classical distinguished triangle]
  \label{thm:external-synthetic-triangle-lift}
  \uses{thm:external-nu-cofiber-criterion,thm:external-synthetic-lift,def:synthetic-lambda-quotient,def:external-result}
  \notready
  For a distinguished triangle
  $X\xrightarrow fY\xrightarrow gZ\xrightarrow h\Sigma X$ with
  $\operatorname{AF}(h)>0$ and short exact $H\mathbb F_2$-homology, an
  explicit BHS input supplies a synthetic distinguished triangle ending in
  $\Sigma^{1,0}\nu X$ and a particular triangle-compatible map $\widehat h$
  satisfying $\nu h=\lambda\widehat h$.  If
  $\operatorname{AF}(h)=k<\infty$ and a full lift $\widetilde h$ is chosen
  separately, the source identifies $\widehat h$ with
  $\lambda^{k-1}\widetilde h$ only modulo a $\lambda$-torsion map, not by an
  equality of maps.
\end{theorem}

\begin{proof}
  Use the cofiber criterion on the first two maps and retain the located
  construction in the proof of BHS Lemma 9.15 as the chosen connecting map.
  Verify the suspension
  $\Sigma^{0,-1}\nu\Sigma X=\Sigma^{1,0}\nu X$ in AIM grading.  The final
  torsion qualification is exactly the comparison recorded in
  \texttt{aimpaper/main.tex:942--943}.
\end{proof}

% SOURCE: BHS Proposition A.13; aimpaper/main.tex:1016-1019.
\begin{theorem}[$\lambda$-adic completeness]
  \label{thm:external-synthetic-lambda-complete}
  \uses{def:lambda-rho-delta-triangle,def:external-result}
  \notready
  An explicit BHS external result supplies a natural equivalence
  \[
    \nu X\simeq\varprojlim_m\nu X/\lambda^m
  \]
  compatible with the quotient maps and the finite
  $\lambda$--$\rho$--$\delta$ triangles.
\end{theorem}

\begin{proof}
  Instantiate BHS Proposition A.13 and transport its grading and inverse-system
  maps to the conventions of Definition~\ref{def:lambda-rho-delta-triangle}.
\end{proof}

\section{Published Kervaire and geometric results}

% SOURCE NOTE: Barratt--Jones--Mahowald inductive method and Burklund--Xu
% Proposition 7.19; aimpaper/main.tex:2117--2125, whose quantifier is ``for
% some theta_5''.  The Burklund--Xu proof internally uses the commutative
% algebra tower of Construction 7.7, but that proof dependency is represented
% only by this external-result source note and does not enter the project graph.
\begin{theorem}[BJM/BX criterion for the source choice]
  \label{thm:external-bjm-bx-criterion}
  \uses{def:synthetic-theta5-choice,def:external-result}
  \notready
  An explicit external-result record supplies a distinguished order-two class
  $\theta_5^{\mathrm{src}}$ detected by $h_5^2$ for which:
  \begin{enumerate}
    \item for every integer $r\ge1$, $h_6^2$ survives to $E_{r+3}$ if and
      only if $\lambda\eta(\theta_5^{\mathrm{src}})^2=0$ in
      $\pi_{125,129}(S^{0,0}/\lambda^{r+1})$;
    \item $h_6^2$ is permanent if and only if
      $\lambda\eta(\theta_5^{\mathrm{src}})^2=0$ in
      $\pi_{125,129}S^{0,0}$.
  \end{enumerate}
  The witness and the two equivalences belong to the same source-bearing
  record, which also includes the grading translation from Burklund--Xu to AIM.
\end{theorem}

\begin{proof}
  This node is the located external input, not a project proof.  Retain the
  distinguished class used by the BJM/BX construction, and apply the explicit
  suspension and grading adapter to its finite and infinite statements.  No
  assertion for an arbitrary choice of $\theta_5$ is included here.
\end{proof}

% SOURCE: aimpaper/main.tex:2127-2139, Remarks 7.4-7.5; Xu and IWX inputs.
\begin{theorem}[External $\theta_5$ order and choice comparison]
  \label{thm:external-theta5-order-data}
  \uses{def:synthetic-theta5-choice,def:external-result}
  \notready
  A named Xu/IWX external-result value supplies that every classical
  $\theta_5$ detected by $h_5^2$ has order two and that the difference of any
  two choices has Adams filtration at least six.  The value carries exact
  theorem locators and the comparison from the source's classical convention
  to Definition~\ref{def:synthetic-theta5-choice}.
\end{theorem}

\begin{proof}
  This is a pure literature input.  Store the two quantified conclusions and
  their grading adapter in one named field; do not derive either conclusion
  from the near-$126$ tables.
\end{proof}

% SOURCE: aimpaper/main.tex:2133-2139, Remark rem:theta5choice, applying the construction in the proof of Burklund--Xu Proposition 7.19 to an order-two bottom-cell choice.
\begin{theorem}[Total differential identity for an order-two choice]
  \label{thm:external-total-differential-identity}
  \uses{def:synthetic-theta5-choice,def:lambda-rho-delta-triangle,def:external-result}
  \notready
  For every order-two synthetic choice $\theta_5$, the located quadratic-
  construction input supplies the identity
  \[
    \delta_1(h_6^2)=\lambda\eta\theta_5^2
  \]
  in the $\delta_1:S^{0,0}/\lambda\to S^{1,-1}$ total-extension problem, with
  its representative and grading conventions.
\end{theorem}

\begin{proof}
  This is the external total-differential comparison asserted in the cited AIM
  remark: use the chosen order-two class as the bottom-cell restriction in the
  quadratic construction, then translate the target suspension to $S^{1,-1}$.
  The project does not reconstruct that construction internally.
\end{proof}

% SOURCE: Browder, The Kervaire invariant of framed manifolds and its generalization; aimpaper/main.tex:117-123, Theorem thm:browder.
\begin{theorem}[Browder criterion]
  \label{thm:external-browder-criterion}
  \uses{def:framed-kervaire-context,def:classical-adams-ss,def:adams-hj,def:ss-permanent-einfty-class,def:external-result}
  \notready
  Fix a stable context $\mathcal C$, a framed Kervaire context over
  $\mathcal C$, and a chosen classical Adams spectral sequence
  $A_{\mathrm{cl}}$ of $S^0_{\mathcal C}$.  An explicit Browder external result
  supplies that a framed manifold with Kervaire invariant one exists in
  dimension $n$ if and only if $n=2^{j+1}-2$ for some $j\ge1$ and the class
  $h_j^2$ is a permanent cycle in this same $A_{\mathrm{cl}}$.  The record must
  carry the precise theorem or page locator in Browder's paper and the
  Pontryagin--Thom convention; it cannot quantify over or silently select a
  second sphere Adams spectral sequence.
\end{theorem}

\begin{proof}
  Instantiate the located Browder theorem and transport its Adams grading and
  framed-cobordism notation into the two project interfaces.  This node remains
  an explicit literature parameter; no project proof of Browder's criterion is
  claimed.
\end{proof}

% SOURCE: aimpaper/main.tex:151-162 and cited classical Kervaire literature.
\begin{theorem}[Prior low-dimensional Kervaire existence]
  \label{thm:external-low-kervaire-existence}
  \uses{def:framed-kervaire-context,def:external-result}
  \notready
  Explicit literature results supply framed Kervaire-invariant-one manifolds
  in dimensions $2,6,14,30,$ and $62$, with one source locator for each
  dimension or a single theorem covering the complete list.
\end{theorem}

\begin{proof}
  Assemble the five located external results into a finite family indexed by
  the stated dimensions, verifying that each conclusion uses the same framed
  and Kervaire conventions as the project context.
\end{proof}

% SOURCE: Hill--Hopkins--Ravenel; cited at aimpaper/main.tex:86-92 and 151-160.
\begin{theorem}[HHR nonexistence beyond dimension $126$]
  \label{thm:external-hhr-nonexistence}
  \uses{def:framed-kervaire-context,def:external-result}
  \notready
  An explicit Hill--Hopkins--Ravenel result supplies nonexistence of
  Kervaire-invariant-one framed manifolds in dimensions
  $2^{j+1}-2$ for $j\ge7$, in the precise form and range of the cited theorem.
\end{theorem}

\begin{proof}
  Import the located HHR statement as an external-result parameter and align
  its dimension indexing and geometric conventions with the project context.
\end{proof}

% SOURCE: Barratt--Jones--Mahowald induction; aimpaper/main.tex:162.
\begin{theorem}[BJM induction on order-two Kervaire classes]
  \label{thm:external-bjm-induction}
  \uses{def:adams-detection,def:external-result}
  \notready
  An explicit Barratt--Jones--Mahowald result supplies: if a class $\theta_j$
  detected by $h_j^2$ satisfies $2\theta_j=0$ and $\theta_j^2=0$, then there is
  an order-two class $\theta_{j+1}$ detected by $h_{j+1}^2$.  This result
  motivates the two open propositions but is not used to prove
  $h_6^2$ permanent.
\end{theorem}

\begin{proof}
  Store the exact inductive theorem and its range as a source-bearing external
  result.  The Blueprint records no converse and does not turn either open
  hypothesis into an assumption.
\end{proof}
